\documentclass[10pt]{beamer}
\usetheme{metropolis}
\usepackage{graphicx}
\usepackage{listings}
\usepackage{booktabs}
\usepackage{xcolor}
\usepackage{hyperref}
\usepackage{tikz}

% ── Listing styles ──────────────────────────────────────────────
\definecolor{codebg}{HTML}{F5F5F5}
\definecolor{codegreen}{HTML}{2E7D32}
\definecolor{codegray}{HTML}{757575}
\definecolor{codeblue}{HTML}{1565C0}

\lstdefinestyle{rstyle}{
  language=R, backgroundcolor=\color{codebg},
  basicstyle=\ttfamily\scriptsize,
  keywordstyle=\color{codeblue}\bfseries,
  stringstyle=\color{codegreen},
  commentstyle=\color{codegray}\itshape,
  breaklines=true, frame=single,
  rulecolor=\color{codegray!30}, showstringspaces=false, tabsize=2,
  morekeywords={library,ggplot,aes,geom_point,scale_color_manual,
    scale_size_continuous,theme_minimal,labs,facet_wrap,annotate}
}
\lstdefinestyle{pythonstyle}{
  language=Python, backgroundcolor=\color{codebg},
  basicstyle=\ttfamily\scriptsize,
  keywordstyle=\color{codeblue}\bfseries,
  stringstyle=\color{codegreen},
  commentstyle=\color{codegray}\itshape,
  breaklines=true, frame=single,
  rulecolor=\color{codegray!30}, showstringspaces=false, tabsize=4,
  morekeywords={import,from,as,pd,plt,sns,np,fig,ax}
}

% ── Metadata ────────────────────────────────────────────────────
\title{Visual Perception \& Pre-Attentive Attributes}
\subtitle{Workshop 1 --- Module 3}
\author{Prof.Asoc.\ Endri Raco}
\institute{Data Visualization for Data Scientists \\ UNYT Tirana}
\date{2025/26}

\begin{document}

% ================================================================
\begin{frame}
  \titlepage
\end{frame}

% ================================================================
\begin{frame}{Module Roadmap}
  \begin{enumerate}
    \item Pre-attentive visual attributes
    \item Feature search vs.\ conjunction search
    \item The Gestalt laws of perceptual grouping
    \item Visual encoding channels and their effectiveness
    \item Weber's law and just-noticeable differences
  \end{enumerate}
  \begin{exampleblock}{Learning Outcome}
    You will be able to select visual encoding channels that exploit the
    human visual system rather than fighting against it.
  \end{exampleblock}
\end{frame}

% ================================================================
\section{Pre-Attentive Processing}
% ================================================================

\begin{frame}{What Is Pre-Attentive Processing?}
  \begin{columns}[T]
    \begin{column}{0.55\textwidth}
      The visual cortex detects certain properties
      \textbf{before conscious attention} engages:
      \begin{itemize}
        \item Processing time: $<$ 200--250 ms
        \item Independent of the number of distractors
        \item Driven by the \emph{primary visual cortex} (V1)
        \item Requires a single distinguishing feature
      \end{itemize}
    \end{column}
    \begin{column}{0.42\textwidth}
      \textbf{Four canonical channels:}
      \begin{enumerate}
        \item Color (hue)
        \item Shape (form)
        \item Size (length / area)
        \item Orientation (angle)
      \end{enumerate}
    \end{column}
  \end{columns}

  \vspace{4pt}
  \emph{Source: Healey, C.~G. \& Enns, J.~T. (2012). Attention and visual
  memory in visualization and computer graphics. IEEE TVCG, 18(7).}
\end{frame}

% ================================================================
\begin{frame}{Pre-Attentive Pop-Out: Color}
  \begin{center}
    \includegraphics[width=0.92\textwidth]{figures_m03/preattentive_color}
  \end{center}
  \begin{alertblock}{Pro-Tip}
    Use colour sparingly: a single accent hue against a neutral field
    produces instant pop-out. Using $>$ 5 hues destroys the effect.
  \end{alertblock}
\end{frame}

% ================================================================
\begin{frame}{Pre-Attentive Pop-Out: Shape}
  \begin{center}
    \includegraphics[width=0.92\textwidth]{figures_m03/preattentive_shape}
  \end{center}
  \begin{exampleblock}{Data Insight}
    Shape pop-out works only for simple, high-contrast distinctions
    (circle vs.\ square). Subtle shape differences (pentagon vs.\ hexagon)
    require serial scanning.
  \end{exampleblock}
\end{frame}

% ================================================================
\begin{frame}{Pre-Attentive Pop-Out: Size \& Orientation}
  \begin{columns}[T]
    \begin{column}{0.48\textwidth}
      \includegraphics[width=\textwidth]{figures_m03/preattentive_size}
    \end{column}
    \begin{column}{0.48\textwidth}
      \includegraphics[width=\textwidth]{figures_m03/preattentive_orientation}
    \end{column}
  \end{columns}
  \vspace{4pt}
  Both size and orientation support pre-attentive detection when the
  difference is sufficiently large. \textbf{Weber's law} (later in this
  module) quantifies ``sufficiently large.''
\end{frame}

% ================================================================
\begin{frame}{Conjunction Search: When Pop-Out Fails}
  \begin{center}
    \includegraphics[width=0.65\textwidth]{figures_m03/conjunction}
  \end{center}
  \begin{alertblock}{Pro-Tip}
    When the target is defined by a \emph{combination} of features (red
    \textbf{AND} square), the visual system must search serially at
    $\sim$50 ms per item. Avoid encoding critical distinctions as
    conjunctions --- use a single unique channel instead.
  \end{alertblock}
\end{frame}

% ================================================================
\begin{frame}{Implications for Chart Design}
  \begin{columns}[T]
    \begin{column}{0.48\textwidth}
      \textbf{Do:}
      \begin{itemize}
        \item Encode the most important distinction using a single
              pre-attentive channel (colour or position)
        \item Limit categorical colours to 5--7 hues
        \item Use size only for ordinal/quantitative data
      \end{itemize}
    \end{column}
    \begin{column}{0.48\textwidth}
      \textbf{Don't:}
      \begin{itemize}
        \item Require viewers to distinguish targets by conjunctions
              (e.g., ``find the red dashed line'')
        \item Use shape alone for $>$ 4 categories
        \item Vary orientation for nominal data
      \end{itemize}
    \end{column}
  \end{columns}
\end{frame}

% ================================================================
\section{Gestalt Laws of Perceptual Grouping}
% ================================================================

\begin{frame}{Why Gestalt Matters for Visualization}
  The Gestalt psychologists (Wertheimer, 1923) identified principles
  by which the visual system organises stimuli into groups:
  \begin{enumerate}
    \item \textbf{Proximity} --- near elements form groups
    \item \textbf{Similarity} --- like elements form groups
    \item \textbf{Continuity} --- smooth paths preferred
    \item \textbf{Closure} --- brain completes incomplete forms
    \item \textbf{Enclosure} --- bounded regions create groups
    \item \textbf{Connection} --- linked elements seen as related
  \end{enumerate}
  \vspace{4pt}
  Every chart layout implicitly invokes these laws.
  Effective design \emph{harnesses} them; poor design \emph{violates} them.
\end{frame}

% ================================================================
\begin{frame}{Proximity}
  \begin{center}
    \includegraphics[width=0.92\textwidth]{figures_m03/gestalt_proximity}
  \end{center}
  \begin{exampleblock}{Data Insight}
    Proximity is the strongest grouping cue. In bar charts, spacing
    between groups must exceed spacing within groups, or the viewer
    will misread the category structure.
  \end{exampleblock}
\end{frame}

% ================================================================
\begin{frame}{Similarity}
  \begin{center}
    \includegraphics[width=0.6\textwidth]{figures_m03/gestalt_similarity}
  \end{center}
  \begin{exampleblock}{Data Insight}
    Despite equal spacing, the alternating colour/shape rows are
    perceived as distinct horizontal groups. This is why colour legends
    work: similar hues are read as ``the same category.''
  \end{exampleblock}
\end{frame}

% ================================================================
\begin{frame}{Continuity \& Closure}
  \begin{center}
    \includegraphics[width=0.92\textwidth]{figures_m03/gestalt_continuity}
  \end{center}
  \begin{itemize}
    \item \textbf{Continuity}: two intersecting sine waves are perceived as
          two smooth paths, not four disjointed segments
    \item \textbf{Closure}: the brain fills in missing arcs to perceive
          complete circles
  \end{itemize}
\end{frame}

% ================================================================
\begin{frame}{Enclosure \& Connection}
  \begin{center}
    \includegraphics[width=0.92\textwidth]{figures_m03/gestalt_enclosure}
  \end{center}
  \begin{alertblock}{Pro-Tip}
    Use enclosure (shaded bands) to highlight groups in scatter plots.
    Use connection (lines between points) to indicate temporal sequence
    or paired comparisons. Both override proximity when present.
  \end{alertblock}
\end{frame}

% ================================================================
\begin{frame}{Gestalt Summary Card}
  \begin{center}
    \includegraphics[width=0.88\textwidth]{figures_m03/gestalt_summary}
  \end{center}
\end{frame}

% ================================================================
\section{Visual Encoding Channels}
% ================================================================

\begin{frame}{Cleveland \& McGill's Channel Ranking}
  \begin{center}
    \includegraphics[width=0.82\textwidth]{figures_m03/channel_ranking}
  \end{center}
  \begin{exampleblock}{Data Insight}
    Cleveland \& McGill (1984) showed experimentally that position
    judgments are $\sim$2$\times$ more accurate than length judgments
    and $\sim$5$\times$ more accurate than area judgments.
  \end{exampleblock}
\end{frame}

% ================================================================
\begin{frame}{Weber's Law: Just-Noticeable Difference}
  \begin{center}
    \includegraphics[width=0.92\textwidth]{figures_m03/weber}
  \end{center}
  \vspace{2pt}
  \textbf{Weber's law}: $\Delta S / S = k$ (constant). The smallest
  detectable change is a fixed \emph{proportion} of the stimulus, not a
  fixed amount. Implication: small differences in area or colour
  saturation are invisible; position differences are always detectable.
\end{frame}

% ================================================================
\begin{frame}{Change Blindness}
  \begin{center}
    \includegraphics[width=0.92\textwidth]{figures_m03/change_blindness}
  \end{center}
  \begin{alertblock}{Pro-Tip}
    When comparing two views (before/after, Version A/B), viewers miss
    subtle changes without direct annotation. Always highlight the
    specific data points that changed using colour, annotation, or
    connecting lines.
  \end{alertblock}
\end{frame}

% ================================================================
\begin{frame}{Choosing Channels: A Decision Framework}
  \centering
  \begin{tabular}{lll}
    \toprule
    \textbf{Data Type} & \textbf{Best Channel} & \textbf{Avoid} \\
    \midrule
    Quantitative   & Position, length       & Area, colour saturation \\
    Ordinal        & Position, size, saturation & Shape \\
    Nominal ($\leq$7) & Colour hue, shape    & Position (wastes axis) \\
    Nominal ($>$7)    & Faceting (small mult.)  & Colour (too many hues) \\
    \bottomrule
  \end{tabular}

  \vspace{8pt}
  \begin{exampleblock}{Data Insight}
    Map the most important variable to the most effective channel.
    Secondary variables get secondary channels. Never encode the same
    variable redundantly on two channels unless deliberately reinforcing.
  \end{exampleblock}
\end{frame}

% ================================================================
\section{Encoding Channels in Code}
% ================================================================

\begin{frame}[fragile]{Multi-Channel Encoding in R}
  \begin{columns}[T]
    \begin{column}{0.52\textwidth}
\begin{lstlisting}[style=rstyle]
library(ggplot2)

# Four simultaneous channels:
# x = position (spend)
# y = position (ROI)
# color = category (nominal)
# size = budget (quantitative)

ggplot(df, aes(
    x = spend,
    y = roi,
    color = category,
    size = budget)) +
  geom_point(alpha = 0.6) +
  scale_color_manual(
    values = c(A = "#1565C0",
               B = "#E53935",
               C = "#2E7D32")) +
  scale_size_continuous(
    range = c(2, 10)) +
  theme_minimal() +
  labs(title = "4-Channel Plot")
\end{lstlisting}
    \end{column}
    \begin{column}{0.45\textwidth}
      \includegraphics[width=\textwidth]{figures_m03/encoding_r}
    \end{column}
  \end{columns}
\end{frame}

% ================================================================
\begin{frame}[fragile]{Multi-Channel Encoding in Python}
  \begin{columns}[T]
    \begin{column}{0.52\textwidth}
\begin{lstlisting}[style=pythonstyle]
import seaborn as sns
import matplotlib.pyplot as plt

# seaborn handles channels via
# hue, size, style parameters

sns.scatterplot(
    data=df,
    x="spend",
    y="roi",
    hue="category",   # color
    size="budget",     # size
    alpha=0.6,
    palette={
        "A": "#1f77b4",
        "B": "#d62728",
        "C": "#2ca02c"},
    sizes=(30, 250),
)

plt.title("4-Channel Plot")
plt.tight_layout()
\end{lstlisting}
    \end{column}
    \begin{column}{0.45\textwidth}
      \includegraphics[width=\textwidth]{figures_m03/encoding_py}
    \end{column}
  \end{columns}
\end{frame}

% ================================================================
\begin{frame}[fragile]{Gestalt-Aware Layout in R}
  \begin{columns}[T]
    \begin{column}{0.52\textwidth}
\begin{lstlisting}[style=rstyle]
# Proximity: use faceting to
# separate groups spatially
ggplot(df, aes(x = x, y = y)) +
  geom_point(aes(color = group),
             size = 2, alpha = 0.6) +
  facet_wrap(~group, ncol = 3) +
  theme_minimal() +
  theme(legend.position = "none")

# Enclosure: add shaded band
ggplot(df, aes(x = date, y = val)) +
  annotate("rect",
    xmin = as.Date("2023-03-01"),
    xmax = as.Date("2023-06-01"),
    ymin = -Inf, ymax = Inf,
    fill = "#E3F2FD", alpha = 0.5) +
  geom_line(color = "#1565C0") +
  theme_minimal()
\end{lstlisting}
    \end{column}
    \begin{column}{0.45\textwidth}
      \textbf{Design patterns:}

      \begin{itemize}
        \item \textbf{Proximity} $\rightarrow$ faceting
              separates groups physically
        \item \textbf{Enclosure} $\rightarrow$ \texttt{annotate("rect")}
              highlights a time period
        \item \textbf{Connection} $\rightarrow$ \texttt{geom\_line}
              links sequential points
        \item \textbf{Similarity} $\rightarrow$ consistent
              colour palette
      \end{itemize}
    \end{column}
  \end{columns}
\end{frame}

% ================================================================
\begin{frame}[fragile]{Gestalt-Aware Layout in Python}
  \begin{columns}[T]
    \begin{column}{0.52\textwidth}
\begin{lstlisting}[style=pythonstyle]
import matplotlib.pyplot as plt

fig, axes = plt.subplots(1, 3,
    sharey=True, figsize=(10, 3))

# Proximity via subplots
for ax, grp in zip(axes, groups):
    sub = df[df["group"] == grp]
    ax.scatter(sub["x"], sub["y"],
        c="#1565C0", s=25, alpha=0.6)
    ax.set_title(grp)

# Enclosure via axvspan
fig, ax = plt.subplots()
ax.axvspan("2023-03", "2023-06",
    color="#E3F2FD", alpha=0.5)
ax.plot(dates, values,
    color="#1565C0")

# Connection via plot()
ax.plot(x, y, "o-",
    color="#E53935", lw=1.5)
\end{lstlisting}
    \end{column}
    \begin{column}{0.45\textwidth}
      \textbf{Python equivalents:}

      \begin{itemize}
        \item \textbf{Proximity} $\rightarrow$ \texttt{subplots}
              with shared axes
        \item \textbf{Enclosure} $\rightarrow$ \texttt{ax.axvspan()}
              for shaded bands
        \item \textbf{Connection} $\rightarrow$ \texttt{plot("o-")}
              links markers
        \item \textbf{Similarity} $\rightarrow$ palette dict
              for consistent colours
      \end{itemize}
    \end{column}
  \end{columns}
\end{frame}

% ================================================================
\section{Synthesis}
% ================================================================

\begin{frame}{Key Takeaways}
  \begin{enumerate}
    \item \textbf{Pre-attentive attributes} (colour, shape, size,
          orientation) are detected in $<$ 200 ms --- use them for the
          most critical data distinction.
    \item \textbf{Conjunction search} is slow ($\sim$50 ms/item) ---
          never require viewers to combine two features to find a target.
    \item \textbf{Gestalt laws} (proximity, similarity, continuity,
          closure, enclosure, connection) govern how elements are grouped
          --- harness them or suffer misreadings.
    \item \textbf{Channel effectiveness} (Cleveland \& McGill, 1984):
          position $>$ length $>$ angle $>$ area $>$ colour --- map the
          most important variable to the best channel.
    \item \textbf{Weber's law}: the just-noticeable difference is
          proportional to the stimulus. Small data differences require
          high-precision channels (position), not low-precision ones
          (area, saturation).
  \end{enumerate}
\end{frame}

% ================================================================
\begin{frame}{Essential Readings for This Module}
  \begin{itemize}
    \item Cleveland, W.~S. \& McGill, R. (1984). ``Graphical Perception:
          Theory, Experimentation, and Application.'' \emph{JASA}, 79(387).
    \item Healey, C.~G. \& Enns, J.~T. (2012). ``Attention and Visual
          Memory in Visualization.'' \emph{IEEE TVCG}, 18(7).
    \item Ware, C. (2021). \emph{Information Visualization: Perception
          for Design}, 4th ed., Chapters 1--5.
    \item Few, S. (2012). \emph{Show Me the Numbers}, Chapters 4--5.
    \item Munzner, T. (2014). \emph{Visualization Analysis and Design},
          Chapter 5 (Marks and Channels).
  \end{itemize}

  \vspace{6pt}
  \textbf{Next Module}: Color Theory \& Accessibility (W01-M04)
\end{frame}

\end{document}
